\subsection*{15.3--15.6　初等奇点的局部刻画}
\begingroup
\def\OO{\mathcal O}\def\m{\mathfrak m}\def\mm{\mathfrak m}\def\Sym{\operatorname{Sym}}
\noindent\textbf{15.3.}\quad 将上文（SUP）命题8“推广”到重数$n$的情形。\footnote{法文清楚写作 SUP，即上文；英译本 15.3 排作 $X\cup P$，15.4 也把 SUP 排为 $S\cup P$，此处不当成数学符号。}

\noindent\textbf{15.4.}\quad 将上文（SUP）命题9和推论10“推广”到重数$n$的情形。（这其实不是推广，而是一些变式……。）

\noindent\textbf{15.5.}\quad 对局部诺特的$X$，引入$\OO_X$的一个截面$\varphi$，或更一般地，$X$上一个局部自由模的截面的\emph{重数$n$的初等奇异零点}的概念（当$n=2$时，为初等二次奇异零点，或称普通二次奇异零点），以及相应的几何概念（底域为$k$）。

\noindent\textbf{15.6.}\quad 设$X$是局部诺特正则预概形$X'$的子预概形，$x\in X$。$x$是$X$的重数$n$的初等奇点，当且仅当存在$x$在$X'$中的一个开邻域$U$及一条$\OO_U$-正则序列$(\varphi_1,\ldots,\varphi_d)$，使
\[
X|_U=V(\varphi_1,\ldots,\varphi_d),
\]
且$x$是$(\varphi_1,\ldots,\varphi_d)$的重数$n$的初等奇异零点。
特别地，若$A$是重数$n$的诺特初等局部环，则$A$是“完全交”环。再者，借助15.4可看出，在$x$的一个邻域内可以找到$X'$的一个包含$X$的正则子预概形$X''$；在$X''$中可用单一方程$\varphi=0$描述$X$，而$x$是$\varphi$的一个重数$n$的初等奇异零点。\footnote{法文物理159页在“$X''$是……的子预概形”处印的是$X$，漏掉撇号；英译作“of $x$”。按“$X''$包含$X$”及前文的正则环境，正文恢复为$X\subseteq X''\subseteq X'$。}

特别可知，若$A$为重数$n$的诺特初等局部环，则$A$是 Cohen--Macaulay 环。我们将在第20节证明（用一个简单的吹起计算），$\Spec A$的闭点是它唯一的奇点。于是，$A$正规当且仅当$\dim A\ge2$，约化当且仅当$\dim A\ge1$。若$\dim A=0$，则$A$为重数$n$的初等环，当且仅当$\m/\m^2$的秩为$1$，且$n$是使$\m^n=0$的最小整数；注意，这些环也是离散赋值环模其极大理想$n$次幂所得的商。
\paragraph{译校说明（非原文）。}
\paragraph{秩估计须以关系为前提。}
§15.3只是提出重数$n$的变式，不是原文已给出的命题。一种局部域上的核验是：设$P=\Sym_F(E)$，$q\in P_n$为已有的非零关系，目标$G$是由$P$满射得到的分次环。乘$q$在$P$中为单射，故
\[
\dim_F(P/(q))_j=\dim_F P_j-\dim_F P_{j-n},\quad P_i=0\ (i<0).
\]
若$\dim_F G_j$又以上式右端为下界，则比较逐次满射$(P/(q))_j\twoheadrightarrow G_j$的维数，得到各次同构。若不先有关系，只有这些下界并不够：$G=F[T_1,T_2]$，由它自身作源的分次同态满足所有这样的下界，却无非零关系。相对情形下须另检查关系线及各次的局部自由性等假设。§15.4和§15.5同样只保留拟稿任务，未补作原文结论或定义。

\paragraph{从正则环境删去线性方程。}
在$x$处写$A=R/I$，$R$是正则局部环。选$I$中元素，其线性类构成切余空间映射
\[
\m_R/\m_R^2\longrightarrow\m_A/\m_A^2
\]
的核的一组基。它们可延成$R$的一组正则参数。商去这些元素得到正则局部环$B$，设余下理想为$J$，使$A=B/J$。$B$与$A$的切余空间相同，因而由\href{https://stacks.math.columbia.edu/tag/00NO}{正则局部环的分次环定理}和§15.2，$\gr_{\m_B}(B)\to\gr_{\m_A}(A)$的核是初始理想$(q)$。取$f\in J$，使其初始型为$q$。对$h\in J$，其非零初始项可由$f$的倍数抵消；逐次重复，在每个阶数$N$上得到$h\in(f)+\m_B^N$。在$B/(f)$应用\href{https://stacks.math.columbia.edu/tag/00IP}{Krull 交定理}，即有$h\in(f)$，所以$J=(f)$。反过来，分次多项式环是整环，对单方程的倍数初始型相乘，所以一个具有光滑次数$n$初始型的方程给出初始理想$(q)$和§15.2的结构。线性参数及其后非零的$f$组成正则序列。这是嵌入情形在$x$处的局部环补证，不是§15.5原本未写出的普遍定义。这个论证本身也不借助任意诺特概形的正则轨迹为开这一假设。

对于原文任意$A$为完全交的断语，采用诺特局部环的标准意义，即其完备化是正则局部环除以正则序列所得。完备化不改变极大理想分次环；由\href{https://stacks.math.columbia.edu/tag/032A}{Cohen 结构定理}取正则局部环到该完备化的满射，再作上述约化。已给定正则环境的情形则无须先取完备化。
\paragraph{吹起与零维情形。}
\paragraph{闭点之外正则。}
设$A$为重数$n\ge2$的诺特初等局部环。将$\Spec A$沿极大理想$\m$吹起。吹起在闭点之外是同构，例外纤维为$\Proj\gr_\m(A)$，也是\href{https://stacks.math.columbia.edu/tag/01OF}{有效 Cartier 除子}。由初始型光滑，该纤维正则。因此在纤维上每一点，吹起的局部环模一个非零因子所得的是正则局部环，\href{https://stacks.math.columbia.edu/tag/00NU}{原局部环也正则}。吹起为固有态射，每一点的闭包在底的像都是含闭点的非空闭集，故该点总有一个落在闭纤维上的特化；再以正则局部环的局部化仍正则，得吹起处处正则。故$\Spec A$的穿孔谱正则。零维时穿孔谱为空。维数至少一时，Cohen--Macaulay 性给出$S_1$，穿孔谱的正则性给出$R_0$，因而约化；维数至少二时用$S_2$和$R_1$得到正规。闭点非正则 (§15.2)，故$A$不是域；在维数零时不约化，而在维数零或一时若正规，则分别为域或离散赋值环，均矛盾。

\paragraph{零维的离散赋值环表示。}
在维数零，§15.2的分次环形状及维数公式迫使嵌入维数为一。反之，若嵌入维数为一，由 Nakayama 引理有$\m=(\pi)$；若$n$为最小消失指数，则各次至$n-1$次的商皆为一维，故$\gr_\m(A)=F[T]/(T^n)$，其中$F=A/\m$。次数$n$型在射影空间$\mathbf P^0_F$截出的概形为空，因而光滑。这是\emph{分次环}的表达，不把混合特征的$A$误认成$F[T]/(T^n)$。

同特征时，Cohen 结构定理给出系数域，取$F[[T]]\to A$，$T\mapsto\pi$。混合特征且$F$的特征为$p$时，取 Cohen 系数离散赋值环$C$。若$p\ne0$于$A$，则$p=\pi^eu$，其中$u$是单位、$1\le e<n$。若$e=1$，$C\to A$已在剩余域上满射，且$p$生成$\m$，故按$\m$的有限过滤逐次提升可知该映射满射。若$e\ge2$，先得满射$C[[T]]\to A$（$T\mapsto\pi$），取$u$的提升$U(T)$；$p-T^eU(T)$在线性参数中有非零类$p$，商去它得到一个一维正则局部环，即离散赋值环，仍满射至$A$。以上各满射的核都是源的极大理想$n$次幂，因为源是离散赋值环，且$n$为极大理想的最小消失指数。
\endgroup
